% Figure: the tromino induction --- base, the quartering step, and the result.
% Source for Figure~\ref{fig:tromino}. TikZ only; regenerable by compiling main.tex.
\begin{tikzpicture}[
  cell/.style={draw=rulegray,minimum size=4.6mm,inner sep=0pt},
  hole/.style={fill=amberbg,draw=amber},
  ctr/.style={fill=accentsoft,draw=accent},
  lab/.style={font=\footnotesize\sffamily,text=slate},
]
% ---------- panel 1: base 2x2 ----------
\begin{scope}[xshift=0mm]
\foreach \p/\v in {0/0/{},0/1/B,1/0/B,1/1/B}
  \node[cell] at ({mod(\p,2)*0.46},{-int(\p/2)*0.46}) {\v};
\node[cell,hole] at (0,0) {};
\node[lab,anchor=north] at (0.46,-1.05) {base: $2\times2$};
\end{scope}
% ---------- panel 2: quartered 4x4 with centre tromino ----------
\begin{scope}[xshift=26mm]
\foreach \r in {0,...,3}
  \foreach \c in {0,...,3}
    \node[cell] at (\c*0.46,-\r*0.46) {};
\node[cell,hole] at (0,0) {};
% centre tromino (cells (1,2),(2,1),(2,2) in row,col)
\node[cell,ctr] at (2*0.46,-1*0.46) {};
\node[cell,ctr] at (1*0.46,-2*0.46) {};
\node[cell,ctr] at (2*0.46,-2*0.46) {};
% quarter boundaries
\draw[accent,thick] (0.92,0.23) -- (0.92,-1.15);
\draw[accent,thick] (-0.23,-0.92) -- (1.61,-0.92);
\node[lab,anchor=north,align=center]{cut in quarters,\\ one L at centre};
\end{scope}
% ---------- panel 3: subproblem view ----------
\begin{scope}[xshift=56mm]
\foreach \r in {0,...,1}
  \foreach \c in {0,...,1}
    \node[cell] at (\c*0.46,-\r*0.46) {};
\node[cell,hole] at (0,0) {};
\node[cell,ctr] at (1*0.46,-1*0.46) {};
\node[lab,anchor=north,align=center]{each quadrant:\\ a smaller hole};
\end{scope}
% ---------- panel 4: assembled 4x4 (real tiling) ----------
\begin{scope}[xshift=84mm]
\foreach \rc/\row in {0/{H,B,C,C},1/{B,B,A,C},2/{D,A,A,E},3/{D,D,E,E}}
  \foreach \c [count=\i] in \row
    \node[cell] at (\i*0.46-0.46,-\rc*0.46) {\c};
\node[cell,hole] at (0,0) {};
\node[lab,anchor=north] at (0.69,-1.3) {assembled: 5 trominos};
\end{scope}
\end{tikzpicture}
